\documentclass[12pt,a4paper]{amsart}

\usepackage{amsmath,amssymb,amsthm}
\usepackage{mathtools}
\usepackage{hyperref}
\usepackage{geometry}
\geometry{margin=2.5cm}

% -------------------------------------------------------
% Theorem environments
% -------------------------------------------------------
\newtheorem{theorem}{Theorem}[section]
\newtheorem{lemma}[theorem]{Lemma}
\newtheorem{corollary}[theorem]{Corollary}
\newtheorem{proposition}[theorem]{Proposition}
\theoremstyle{definition}
\newtheorem{definition}[theorem]{Definition}
\newtheorem{remark}[theorem]{Remark}
\newtheorem{conjecture}[theorem]{Conjecture}

% -------------------------------------------------------
% Custom commands
% -------------------------------------------------------
\newcommand{\N}{\mathbb{N}}
\newcommand{\Z}{\mathbb{Z}}
\newcommand{\Q}{\mathbb{Q}}
\newcommand{\R}{\mathbb{R}}
\newcommand{\C}{\mathbb{C}}
\newcommand{\Fp}{\mathbb{F}_p}
\DeclareMathOperator{\ord}{ord}
\DeclareMathOperator{\Sha}{\Sha}

% -------------------------------------------------------
\begin{document}
% -------------------------------------------------------

\title{The Birch--Swinnerton-Dyer Conjecture}

\author{Michael Fuhrmann}
\address{Specialist Mathematics Project, Build 84}
\email{specialist-maths@localhost}
\date{\today}

\begin{abstract}
The Birch--Swinnerton-Dyer (BSD) conjecture is one of the seven Clay
Millennium Prize Problems.
It predicts a deep connection between the arithmetic of an elliptic curve
$E/\Q$ (its rank $r$, torsion, and period lattice) and the analytic
behaviour of its $L$-function $L(E,s)$ at $s=1$.

We state the \emph{weak BSD}: $L(E,1) = 0 \iff E(\Q)$ is infinite.
We state the \emph{strong BSD}: $\ord_{s=1}L(E,s) = r$, and give the
full conjectural formula for the leading coefficient involving
$\Omega_E$, the real period, the Tamagawa numbers $c_p$,
the Tate--Shafarevich group $\Sha(E/\Q)$, and the regulator $R_E$.

We describe the known cases: rank~$0$ and rank~$1$ (Coates--Wiles 1977,
Kolyvagin 1990), and the heuristic numerical evidence via the
Birch--Swinnerton-Dyer product $\prod_{p \le X} \#E(\Fp)/p \sim C(\log X)^r$.
\end{abstract}

\maketitle
\tableofcontents

% ================================================================
\section{Introduction}
% ================================================================

In the 1960s, Birch and Swinnerton-Dyer \cite{BSD1965} computed
$\prod_{p \le X} \#E(\Fp)/p$ for many elliptic curves $E$ and observed
numerically that this product grows like $C \cdot (\log X)^r$ where $r$
is the rank of $E(\Q)$.
This led them to conjecture that the order of vanishing of $L(E,s)$
at $s = 1$ equals the rank.

The conjecture is now split into two parts:
the \emph{weak BSD} (rank = order of vanishing) and the
\emph{strong BSD} (precise formula for the leading coefficient).

% ================================================================
\section{Weak BSD: Rank and the Order of Vanishing}
% ================================================================

\begin{conjecture}[Weak BSD]
\label{conj:weak_bsd}
Let $E$ be an elliptic curve over $\Q$ with rank $r = \mathrm{rank}\, E(\Q)$.
Then
\[
  \ord_{s=1} L(E,s) \;=\; r.
\]
In particular:
\begin{itemize}
  \item $L(E,1) \ne 0 \iff r = 0 \iff E(\Q)$ is finite.
  \item $L(E,1) = 0$ and $L'(E,1) \ne 0 \iff r = 1$.
\end{itemize}
\end{conjecture}

\begin{remark}[Current status]
The weak BSD is proved only in special cases:
\begin{itemize}
  \item Rank $0$: Coates--Wiles (1977) \cite{CoatesWiles1977} proved $r=0$
    when $E$ has complex multiplication and $L(E,1) \ne 0$.
  \item Kolyvagin (1990) \cite{Kolyvagin1990}: For a modular elliptic curve
    with $L(E,1) \ne 0$, $E(\Q)$ is finite and $\Sha(E/\Q)$ is finite.
    For $L'(E,1) \ne 0$ (and sign $-1$), $r = 1$.
  \item Gross--Zagier (1986): Under certain conditions, $L'(E,1) \ne 0$
    implies the existence of a rational point of infinite order.
  \item The general case is open.
\end{itemize}
\end{remark}

% ================================================================
\section{The Tate--Shafarevich Group}
% ================================================================

\begin{definition}[Tate--Shafarevich group]
\label{def:sha}
Let $E/\Q$.  The \emph{Tate--Shafarevich group} is
\[
  \Sha(E/\Q) \;=\; \ker\!\Bigl(H^1(\Q, E) \to \prod_v H^1(\Q_v, E)\Bigr),
\]
where $v$ ranges over all places (primes $p$ and the archimedean place $\infty$),
and $H^1(\Q, E)$ denotes the Galois cohomology.

Intuitively, $\Sha(E/\Q)$ parametrizes torsors under $E$ that have
points over every completion $\Q_v$ but no rational point over $\Q$
--- i.e.\ failures of the Hasse local-global principle for $E$.
\end{definition}

\begin{proposition}[Basic properties of $\Sha$]
\label{prop:sha_basic}
\begin{enumerate}
  \item[\textup{(i)}] $\Sha(E/\Q)$ is a torsion abelian group.
  \item[\textup{(ii)}] $\Sha(E/\Q)$ is equipped with an alternating
    (skew-symmetric) bilinear pairing (the Cassels--Tate pairing),
    so the order $|\Sha(E/\Q)|$ is a perfect square (when finite).
  \item[\textup{(iii)}] Under the assumption of BSD, $\Sha(E/\Q)$ is finite.
  \item[\textup{(iv)}] Kolyvagin proved $\Sha(E/\Q)$ is finite when
    $L(E,1) \ne 0$ or $L'(E,1) \ne 0$ (in the analytic rank $\le 1$ case).
\end{enumerate}
\end{proposition}

% ================================================================
\section{Strong BSD: The Leading Coefficient Formula}
% ================================================================

\begin{conjecture}[Strong BSD]
\label{conj:strong_bsd}
Let $E/\Q$ have rank $r$, real period $\Omega_E$, regulator $R_E$,
Tate--Shafarevich group $\Sha(E/\Q)$ (assumed finite), torsion subgroup
of order $|E(\Q)_{\mathrm{tors}}|$, and Tamagawa numbers $c_p = |E(\Q_p)/E_0(\Q_p)|$
for each prime $p$ of bad reduction.  Then
\begin{equation}
  \label{eq:strong_bsd}
  \lim_{s \to 1} \frac{L(E,s)}{(s-1)^r}
  \;=\;
  \frac{\Omega_E \cdot |\Sha(E/\Q)| \cdot R_E \cdot \prod_p c_p}
       {|E(\Q)_{\mathrm{tors}}|^2}.
\end{equation}
\end{conjecture}

\begin{remark}[The ingredients]
\begin{itemize}
  \item $\Omega_E = \int_{E(\R)} |\omega|$: the real period of the
    N\'eron differential $\omega = dx/2y$.
  \item $R_E = \det(\langle P_i, P_j\rangle)_{1 \le i,j \le r}$: the
    N\'eron--Tate regulator, the determinant of the height pairing matrix
    for a basis $\{P_1,\ldots,P_r\}$ of $E(\Q)/E(\Q)_{\mathrm{tors}}$.
  \item $c_p = |E(\Q_p) / E_0(\Q_p)|$: the Tamagawa number at $p$
    (the number of components of the N\'eron model at $p$).
  \item $\Sha(E/\Q)$: the Tate--Shafarevich group.
  \item $E(\Q)_{\mathrm{tors}}$: the torsion subgroup (Paper~25).
\end{itemize}
The strong BSD formula~\eqref{eq:strong_bsd} has been numerically verified
for many curves but proved only for very special cases.
\end{remark}

% ================================================================
\section{Numerical Evidence: The BSD Product}
% ================================================================

\begin{definition}[BSD heuristic product]
\label{def:bsd_product}
For $X > 0$, define the \emph{BSD heuristic product}
\[
  B_E(X) \;=\; \prod_{p \le X,\; p\text{ good}} \frac{\#E(\Fp)}{p}.
\]
\end{definition}

\begin{theorem}[Birch--Swinnerton-Dyer numerical observation]
\label{thm:bsd_numerical}
Birch and Swinnerton-Dyer \cite{BSD1965} computed $B_E(X)$ for many curves
and observed:
\[
  B_E(X) \;\sim\; C_E\,(\log X)^r \quad (X \to \infty),
\]
where $r = \mathrm{rank}\,E(\Q)$ and $C_E$ is a curve-dependent constant.
\end{theorem}

\begin{proof}[Heuristic justification]
Using Hasse's theorem $\#E(\Fp) = p + 1 - a_p$ with $|a_p| \le 2\sqrt{p}$:
\[
  \frac{\#E(\Fp)}{p} = 1 + \frac{1}{p} - \frac{a_p}{p}
  \approx 1 - \frac{a_p}{p}.
\]
The product $\prod_{p \le X}(1 - a_p/p)$ is related to $L(E,1)$ via
the Euler product (in a formal sense).
When $L(E,1) = 0$ (equivalently rank $r \ge 1$), the product diverges
logarithmically, giving the observed $(\log X)^r$ growth.
This heuristic is made precise by the relation
\[
  -\frac{d}{ds}\log L(E,s)\big|_{s=1} = \sum_p \frac{a_p \log p}{p}
  + \text{lower order},
\]
which controls the growth of partial products.
\end{proof}

% ================================================================
\section{Known Cases}
% ================================================================

\begin{theorem}[Coates--Wiles, 1977]
\label{thm:coates_wiles}
Let $E/\Q$ be an elliptic curve with complex multiplication by an
imaginary quadratic field $K$.  If $L(E,1) \ne 0$, then $E(\Q)$ is finite.
\end{theorem}

\begin{proof}[Proof sketch]
The key tool is the theory of Euler systems for $E$ with complex multiplication.
Coates and Wiles \cite{CoatesWiles1977} used a $p$-adic $L$-function and
the arithmetic of the CM field to show that $L(E,1) \ne 0$ implies the
$p^\infty$-Selmer group $\mathrm{Sel}_{p^\infty}(E/\Q)$ is finite for
all primes $p$ dividing the CM period.  This forces $E(\Q)$ to be finite.
\end{proof}

\begin{theorem}[Kolyvagin, 1990]
\label{thm:kolyvagin}
Let $E/\Q$ be a modular elliptic curve.
\begin{enumerate}
  \item[\textup{(i)}] If $L(E,1) \ne 0$, then $r = 0$ (i.e.\ $E(\Q)$ is finite)
    and $\Sha(E/\Q)$ is finite.
  \item[\textup{(ii)}] If $L'(E,1) \ne 0$ (first derivative nonzero at $s=1$),
    then $r = 1$ and $\Sha(E/\Q)$ is finite.
\end{enumerate}
\end{theorem}

\begin{proof}[Proof sketch]
Kolyvagin's \cite{Kolyvagin1990} proof uses \emph{Euler systems} of Heegner points:
points on $E$ coming from the theory of complex multiplication,
satisfying a norm-compatible system as the field varies.
The key steps:
\begin{enumerate}
  \item Gross--Zagier (1986): For $L'(E,1) \ne 0$, the Heegner point on $E$
    is a point of infinite order.
  \item Kolyvagin's descent: Using Euler system relations, the Selmer group
    has the predicted structure.
  \item The finiteness of $\Sha$ follows from the structure of the
    Kolyvagin system.
\end{enumerate}
\end{proof}

% ================================================================
\section{The Tate--Shafarevich Group and Descent}
% ================================================================

\begin{remark}[Computing the Selmer group]
In practice, one computes the rank by performing a \emph{2-descent}
(or $n$-descent for larger $n$): this gives the \emph{Selmer group}
$\mathrm{Sel}_2(E/\Q)$ fitting in the exact sequence
\[
  0 \to E(\Q)/2E(\Q) \to \mathrm{Sel}_2(E/\Q) \to \Sha(E/\Q)[2] \to 0.
\]
From the Selmer group, one obtains:
\begin{align*}
  r + \dim_{\mathbb{F}_2} \Sha[2]
  &= \dim_{\mathbb{F}_2} \mathrm{Sel}_2(E/\Q) - \dim_{\mathbb{F}_2}E(\Q)_{\mathrm{tors}}[2].
\end{align*}
If $\Sha(E/\Q) = 0$, then $r = \dim_{\mathbb{F}_2} \mathrm{Sel}_2(E/\Q) - \ldots$
gives the exact rank.
\end{remark}

% ================================================================
\section{Conclusion}
% ================================================================

The Birch--Swinnerton-Dyer conjecture predicts a precise link between the
arithmetic of $E(\Q)$ (rank, regulator, Tate--Shafarevich group) and the
analytic behaviour of $L(E,s)$ at $s=1$.

The weak form (rank = analytic rank) is proved only for analytic rank $\le 1$
(Kolyvagin 1990).  The strong form~\eqref{eq:strong_bsd} has been verified
numerically for thousands of curves but remains unproved in general.

BSD is one of the seven Millennium Prize Problems; its resolution would
represent a fundamental advance in our understanding of rational points
on algebraic varieties.

% ================================================================
\begin{thebibliography}{10}

\bibitem{BSD1965}
B.~J.~Birch and H.~P.~F.~Swinnerton-Dyer,
Notes on elliptic curves, II,
\emph{J.\ Reine Angew.\ Math.}\ \textbf{218} (1965), 79--108.

\bibitem{CoatesWiles1977}
J.~Coates and A.~Wiles,
On the conjecture of Birch and Swinnerton-Dyer,
\emph{Invent.\ Math.}\ \textbf{39} (1977), 223--251.

\bibitem{Kolyvagin1990}
V.~A.~Kolyvagin,
Euler systems,
in: \emph{The Grothendieck Festschrift}, Vol.~II,
Progr.\ Math.~87, Birkh\"auser, 1990, pp.~435--483.

\bibitem{GrossZagier1986}
B.~H.~Gross and D.~B.~Zagier,
Heegner points and derivatives of $L$-series,
\emph{Invent.\ Math.}\ \textbf{84} (1986), 225--320.

\bibitem{Silverman1986}
J.~H.~Silverman,
\emph{The Arithmetic of Elliptic Curves},
Graduate Texts in Mathematics~106,
Springer, New York, 1986.

\bibitem{Wiles1995}
A.~Wiles,
Modular elliptic curves and Fermat's last theorem,
\emph{Ann.\ of Math.}\ (2) \textbf{142} (1995), 443--551.

\bibitem{Milne2006}
J.~S.~Milne,
\emph{Elliptic Curves},
BookSurge Publishing, 2006.
Available at \url{https://www.jmilne.org/math/}.

\end{thebibliography}

\end{document}
